\documentclass[../main.tex]{subfiles}
\begin{document}
%==========================================TIÊU ĐỀ TẬP========================================
\begin{table}[h]
  \begin{adjustwidth}{-1cm}{}
    \begin{tabular}{>{\centering\arraybackslash}p{6cm}|>{\raggedright\arraybackslash}p{12.5cm}}
      \multirow{5}{6cm}
      % Cột bên trái: ÉPISODE và số tập
      {\\[-40pt]
        \flaregothic\fontsize{20pt}{20pt}\selectfont \centering
        \color{eptype}ÉPISODE\color{black}\\
        \burbank\fontsize{72pt}{72pt}\selectfont
        \color{epnum}115\color{black}\\[6pt]
        \cafeteria\fontsize{20pt}{20pt}\selectfont\color{epnum}(9-6)\color{black}
        \\[18pt]
        % \flaregothic\fontsize{14pt}{14pt}\selectfont
        % \color{epattr}BIENVENUE, \uline{\color{Banpire} Mel} !
        \color{black}
      }
       &
      % Dòng thứ nhất: tên tập tiếng Nhật
      {\fotskip\fontsize{18pt}{18pt}\selectfont \ruby{七}{たな}\ruby{夕}{ばた}を\ruby{終}{お}わらせない！
      }            \\[6pt]
      % Dòng thứ hai: tên tập tiếng Anh
       & {\cafeteria\fontsize{18pt}{18pt}\selectfont Tanabata Lives On!}                             \\[4pt]
      % Dòng thứ ba: tên tập tiếng Việt
       & {\vnmsans\fontsize{18pt}{18pt}\selectfont (Hậu) Thất tịch}                                  \\[8pt]
      % Dòng thứ tư: Nhân vật xuất hiện
       & {\caviar{\fontsize{14pt}{14pt}\selectfont Nhân vật: \emp{kuon1} \emp{ryuuhou1} \emp{mel} \emp{fubuki} \emp{flare1} \emp{noel} \emp{nene}}}
      \\[8pt]
      % Dòng cuối cùng: Thời gian diễn ra của tập (thường sẽ là ngày phát hành)
       & {\continuum\fontsize{16pt}{16pt}\selectfont Ngày phát hành: 25/07/2021}                     \\[8pt]
       & (dựa theo bản dịch của \href{https://www.facebook.com/mamisubteam}{MamiSubs - Mami Fansub}) \\[18pt]
    \end{tabular}
  \end{adjustwidth}
\end{table}
%=========================================TRUYỆN CHÍNH========================================
\color{black}

\txtbx[2]
{{\color{vol} \uline{Tóm tắt:}} Nene khăng khăng tổ chức Thất tịch giữa tháng Bảy, kéo cả Kuon, Noel, Mel, Flare và Ryuuhou vào cuộc. Nàng Hiệp sĩ vác một cây tre khổng lồ về, rồi cùng nàng Hải cẩu dựng, phá, dựng lại hàng loạt cấu trúc tre kỳ quái: từ dây leo ma quái, trái tim, tàu lượn, cho đến ống mì somen. Sau khi tốn bốn lần sửa, cuối cùng họ cũng dựng được một cây tre đúng nghĩa và treo điều ước. Noel ném cây tre lên trời cho Hikoboshi (vốn dĩ là Fubuki), người đọc điều ước của mọi người: Nene ước mì somen, Noel ước cơm bò, Mel và Flare ước những điều tốt đẹp, Kuon ước nhóm đừng làm hề nữa, còn Ryuuhou ước cho những khoảnh khắc đáng nhớ và ban nhạc của mình. Fubuoshi ban quà cho tất cả mọi người, kết thúc một ngày Thất tịch lệch tận gần ba tuần nhưng đầy ý nghĩa.}

Ngày lễ Thất tịch đã trôi qua được một khoảng thời gian. Dù bầu trời tháng Bảy vẫn còn vương lại chút hơi thở mơ màng của ngày lễ, với phần đông mọi người, khoảnh khắc 7/7 đã là chuyện thuộc về hơn hai tuần trước. Thế nhưng, Momosuzu Nene lại không hề nghĩ theo hướng đó.

Dù lịch đã bước sang ngày Hai Mươi Lăm tháng Bảy, cô nàng Hải cẩu vẫn một mực tin rằng hôm nay chính là\dots\ ngày Thất tịch. Không chút lưỡng lự, cô chạy thẳng tới chiếc bảng trắng trong phòng sinh hoạt, gạch bỏ ngày tháng đang hiển thị rồi nguệch ngoạc viết lên đó dòng chữ to đùng: 「{\fotskip 7月7日}」 hay còn gọi là ``Ngày 7 tháng 7''.

``\textbf{Lễ Thất tịch đã đến rồi!}'' Nene vừa hô to vừa xô đẩy chiếc bảng trắng qua lại tràn đầy hứng khởi, tựa như một nhân vật chính đang mở màn cho lễ hội trường học. Đôi mắt cô long lanh sáng rực, chẳng khác nào những vì sao đang nhấp nháy trên bầu trời đêm.

\img{9-6a}

Đứng bên cạnh, Kuon chỉ biết cười khổ, vừa gãi đầu vừa tỏ ra bất lực. Ngồi kế cậu trên chiếc ghế xoay là Ryuuhou vốn dĩ có thói quen xuất hiện ở bất kỳ buổi tiệc nào mà cậu Tổng được mời, kể cả bữa ăn mừng Thất tịch lần này. Cô bé đang chăm chú theo dõi Nene với gương mặt vừa tò mò thích thú, vừa có chút khó hiểu. Cô chẳng nói gì, chỉ lẳng lặng chuyển ánh mắt sang nhìn anh trai, như thể đang chờ xem màn kịch này sẽ tiếp diễn ra sao.

``Ủa, ngáo hả em? Ngay cả trẻ con mẫu giáo cũng rõ ngày đó đã qua đi nửa tháng rồi mà,'' Kuon thốt lên, giọng đầy sự bất lực, tay còn liên tục xoa trán.

Từ phía sau chiếc bảng trắng, Flare cũng ló đầu ra, thêm vào một câu với vẻ mặt đầy ngờ vực: ``Chịu chưa, ảo tre karate à? Biết rõ là sắp hết tháng Bảy rồi còn gì.''

Ryuuhou cuối cùng cũng mở lời, giọng pha lẫn chút hài hước nhưng vẫn giữ được vẻ điềm nhiên thường có: ``Nene-san, em thấy chị dường như đang sống trong một thế giới riêng của mình quá. Có lẽ chị nên mở lịch trên điện thoại ra xem lại thì hơn đó.''

Nhưng những lời nói của Kuon, Flare và Ryuuhou dường như chẳng thể nào xuyên qua được lớp màn hứng khởi đang bao bọc lấy Nene. Cô nàng vẫn miệt mài chuẩn bị mọi thứ để tạo không khí lễ hội, bất chấp thực tế đã trôi xa ngày Thất tịch gần ba tuần trôi qua.

Cùng lúc đó, Mel bước vào căn phòng với một chiếc hộp nhỏ trên tay, mở nắp ra bên trong là cả bộ dải giấy đủ màu sắc rực rỡ. Cô đưa chúng cho Nene - người đang không ngớt phấn khích - và nói bằng giọng nhẹ nhàng, tựa như một bà tiên đang ban phát những điều ước: ``Dải giấy để viết ước nguyện của em đây này.''

\img{9-6b}

Ryuuhou nhìn những dải giấy đầy màu sắc, đôi mắt cô bé bừng sáng lên vì thích thú nhưng vẫn không quên trêu chọc một câu: ``Nene-san mà cứ thế này thì đến hết năm nay chị cũng chưa kịp viết xong hết mấy điều ước đâu đó.''

Thấy tình hình ấy, Kuon và Flare chỉ biết thở dài, rồi cuối cùng cũng đành xuôi theo dòng chảy\dots\ phi lý của mọi việc.

Nàng Yêu tinh vàng tò mò hỏi thêm: ``Vậy còn cây tre trang trí thì sao? Ai sẽ mang đến?''

Nàng Dơi mỉm cười và đáp một cách rất nhẹ nhàng, giọng vẫn giữ được chút bí ẩn: ``Noel-chan sẽ mang cây tre tới đây thôi\~{}''

Ryuuhou liền nhướng mày nhìn Mel, giọng đầy ngờ vực nhưng vẫn xen lẫn chút tò mò: ``Noel-san mà mang cây tre đến á? Em đâu hay biết chị ấy còn có sở thích làm vườn nữa đấy.''

Trong khi đó, Nene đã gần như dán mắt vào mấy dải giấy đầy màu sắc, ánh mắt cô nàng rực rỡ chẳng khác nào một đứa trẻ sắp được nhận quà lễ hội. Cô lần lượt cầm từng dải giấy lên ngắm nghía, miệng còn lẩm bẩm điều gì đó về những ước nguyện sắp viết lên đó. Kuon đứng khoanh tay nhìn cả nhóm với vẻ cam chịu, nhưng cũng chẳng thể che giấu được nụ cười nửa miệng: ``Thiệt tình là mấy đứa nó vui thật nhỉ\dots''

Ryuuhou quay sang nhìn anh trai, giọng tựa như một nhà bình luận đang đưa ra nhận xét: ``Onii, em thấy các chị ấy vui vẻ là tốt rồi. Ít nhất cũng có chuyện để làm, chứ đừng như bọn mình mà cứ ngồi không chẳng có gì làm.''

Ngay khi cô vừa dứt lời, từ ngoài hành lang vọng vào một tiếng hô lớn: ``\textbf{Hàng tới đây! Hàng tới đây!}''

``À, đến rồi. Để em thử xem trông nó thế nào--'' Flare nói, nhưng vừa mới xoay cái nắm cửa thì\dots\ *RẦM!* Một thứ gì đó dài cồng kềnh đẩy cô nàng Nửa yêu tinh lùi cả mét, cho đến khi cô ngã phịch vào phòng thu âm bên cạnh, tay chân đập loạn xạ lên sàn.

\gif{9-6c}{1}{42}

Ba người còn lại đều choáng váng. Ryuuhou đứng bật dậy khỏi ghế, mắt mở to kinh ngạc. ``C-cái gì thế kia!? Nguyên cả cây tre á!?'' cô bé thốt lên, giọng vừa ngạc nhiên vừa phấn khích.

Noel xuất hiện ngay sau cây tre khổng lồ, vẻ mặt bình thản như thể việc vác nguyên một cái cây từ đâu về là điều hoàn toàn bình thường. Cô đặt cây tre xuống sàn với một tiếng nặng nề, rồi dõng dạc tuyên bố: ``Càng to càng tốt mà!''

``Cây to thế!'' / ``Cái này để vào đâu được trời!?'' Kuon và Mel đồng thanh, vẻ mặt vừa bối rối vừa kinh ngạc, cả hai cùng nhìn cây tre rồi nhìn Noel với ánh mắt không thể tin nổi.

Ryuuhou bước lại gần, đưa tay chạm vào thân tre khổng lồ, giọng đầy thán phục: ``Noel-san, chị lấy cây này ở đâu ra vậy? Trông như vừa chặt từ rừng Amazon về đấy.''

Không đợi ai trả lời, một tiếng nói vang lên từ phía Nene, người vẫn đang đứng với hai ngón tay chỉ lên trời như một thánh nữ của lễ hội: ``Nghệ thuật là vòng lặp của sự sáng tạo và hủy diệt.''

Kuon nuốt khan một cái, mồ hôi túa ra: ``Ý em là\dots\ đập nó ra luôn á!?''

Noel giơ nắm tay ra phía trước, hào hứng nói, cơ bắp cuồn cuộn dưới làn da: ``Cứ để đó cho chị!''

``\dots Thế sao không làm thế từ đầu? Mất công vác nguyên cái cây này làm gì?'' cả Kuon và Mel đều lắc đầu, đưa tay lên trán.

Nene hào hứng giơ hai tay lên cao, như thể đang giải thích một chân lý vũ trụ: ``Cây phải nguyên vẹn thì mới có làm nghệ thuật chứ!''

Ryuuhou nhìn Nene, rồi nhìn cây tre, rồi lại nhìn Noel, giọng vẫn giữ được vẻ hài hước: ``Em nghĩ Nene-san đang có cái nhìn hơi khác về nghệ thuật đấy. Mà thôi, cũng vui.''

``Chị nghĩ không ai để tâm tới mấy cái đó lắm đâu\dots'' Flare thở dài, khoanh tay nhìn chiếc cây tre đang bắt đầu được chặt, giọng đầy bất lực.

Thế là, những người còn lại - trong đó Nene vẫn đang tạo dáng hai ngón tay chỉ lên trời - đứng nhìn như xem xiếc khi Noel bắt đầu ``xử lý'' cây tre. Từng tiếng đập đều đều, chấn động nguyên cả căn phòng, mà họ vẫn không biết cô nàng Hiệp sĩ sẽ làm được gì tiếp theo.

Cô em gái của cậu Tổng đưa tay che miệng, cười khúc khích: ``Em thấy Noel-san đang làm một màn ảo thuật đấy. Cây tre mà bẻ được bằng tay thế này thì chắc cổ tích cũng phải gọi là thánh.''

``Em ấy còn được gọi là `cánh tay đô' của \hll\ mà, mấy cái này nhằm nhò gì,'' ông anh trai đứng khoanh tay nhìn, khẽ tặc lưỡi nhẹ, ``chỉ mỗi tội cái não của em ấy thì chậm phát triển.''

Chỉ với tay không, Noel không mất quá năm phút để bẻ nguyên cây tre thành nhiều đốt đều tăm tắp. Mel, Flare và Nene đều tròn mắt há hốc mồm, trong khi Kuon lặng lẽ đứng khoanh tay, tự nhủ: ``\textit{Nhìn như\dots\ cây tre trăm đốt ấy nhỉ.}'' Ryuuhou thì vỗ tay tán thưởng, giọng đầy hài hước: ``Noel-san, chị mà tham gia Olympic thì chắc không ai địch nổi chị đâu.''

Khi cây tre đã được chia thành những đoạn nhỏ dễ sắp đặt, cả năm người bắt tay vào việc dựng lại ``cây tre ngày lễ Thất tịch nhưng đã quá hạn sử dụng''.

Mel nhún nhảy đầy năng lượng, vỗ tay và hô: ``Yosh, giờ đặt lại thôi!''

Flare và Kuon liếc nhau, nghi ngờ hỏi: ``Có được không đó\dots?''

Nene vẫn không mất một giây do dự, giọng đầy tin tưởng: ``Có Noel-senpai giúp thì chẳng có gì phải lo cả!''

Ryuuhou đứng bên cạnh, nhìn cả nhóm với vẻ hài lòng: ``Onii, em thấy vụ này cũng vui đấy. Mà cây tre này dựng lên xong có khi còn đẹp hơn cả lễ hội thật.''

\img{9-6d}

Kuon khẽ gật đầu, rồi thêm vào như đùa như thật, giọng pha chút hài hước: ``Hoặc là nhờ Luna-chan làm `khắc nhập' với `khắc xuất' cũng được.''

Cả ba người con gái kia quay lại nhìn cậu với ánh mắt bối rối. Flare nhướng mày: ``Cây tre trăm đốt chưa nghe bao giờ à?''

Họ đồng loạt lắc đầu.

``Cũng phải thôi\dots'' cậu thở dài, giọng như đang nhớ về một ký ức xa xăm: ``truyện cổ tích quê anh mà. Mấy đứa không biết cũng phải.''

Ryuuhou nhìn anh trai rồi cười gượng: ``Onii, anh mà đem ra mấy cái điển tích nước ta vào thì họ cũng chẳng hiểu đâu.'' Cô nàng quay sang mọi người rồi chậm rãi giải thích: ``Ở xứ em có một câu chuyện về anh chàng nọ, người mà muốn có một cây tre to lớn để làm nhà. Thay vì chặt trăm cái cây nhỏ, anh ấy lại bẻ một cây tre to thành trăm khúc đều nhau, rồi cuối cùng dùng phép thuật để `khắc nhập' tức gộp chúng lại thành một cây hoàn chỉnh, hay `khắc xuất' để tách chúng ra khi cần. Đó chính là mấy từ Onii nói đấy!''

Nghe xong câu giải thích, Nene mắt tròn xoe, tay vô thức vỗ vào má: ``Waa, phép thuật gộp cây tre á? Thật phi thường luôn! Nếu Noel-senpai biết được cái phép khắc nhập kia, chắc chẳng cần phải tốn sức bẻ cây rồi dựng lại nữa nhỉ?''

Flare cũng nghiêng đầu, gật gù vẻ phấn khích: ``Nghe hay đấy! Tưởng chỉ có chuyện trong truyện cổ của Nhật Bản thôi, hóa ra quê hương của hai anh em cũng có nhiều câu chuyện thú vị tương tự. Cơ mà cổ tích cũng chỉ là cổ tích thôi.''

Mel thì mỉm cười, mái tóc tím khẽ rung khi cô chậm rãi thêm vào: ``Hai câu chuyện có nhiều điểm tương đồng thật nhỉ. Dù rằng chúng ở hai phía khác nhau của thế giới, nhưng người ta lại luôn kể về những điều kì diệu mà sức mạnh của cây tre mang lại.''

Tới lúc này, Noel lại hăng hái giơ cơ bắp lên, giọng đầy tự hào: ``Cơ bắp là câu trả lời cho mọi thứ!''

Ryuuhou nhìn cô nàng Hiệp sĩ, lắc đầu với một nụ cười tinh nghịch: ``Noel-san, em nghĩ cơ bắp của chị đã làm được khá nhiều việc rồi. Nhưng mà không phải lúc nào chị cũng có thể dùng sức mạnh để giải quyết mọi chuyện đâu.''

Kuon lắc đầu, thở ra một câu nhẹ như khói: ``Quan trọng là cái tâm của mình nữa. Không phải cứ đập phá rồi dựng lung tung là được đâu.''

Mel, Flare và Kuon đều nhìn Noel bằng ánh mắt lo ngại pha lẫn bất lực. Duy chỉ có Nene là vẫn sáng mắt lấp lánh, miệng cười toe toét, phụ Noel một tay một cách rất chuyên nghiệp, như thể cả hai là đội kỹ thuật viên chuẩn bị cho một buổi lễ mà chẳng ai biết tên.

Ryuuhou đứng bên cạnh, quan sát cảnh tượng với vẻ mặt đầy thích thú. Cô bé thì thầm với Kuon: ``Onii, em thấy hôm nay sẽ có nhiều chuyện vui đấy.''

\ct

``Bọn em xong rồi đây!'' Noel và Nene cùng reo lên, gương mặt rạng rỡ như vừa hoàn thành một kiệt tác đỉnh cao. Hai người đứng sát bên nhau, tay chống hông đầy vẻ tự hào, chẳng khác nào vừa tạo ra một tác phẩm sắp đặt đáng được trưng bày tại bảo tàng nghệ thuật.

\img{9-6e}

Nhưng thứ nằm trước mắt họ thì\dots\ chẳng ai dám gọi là gì được. Một ``tác phẩm'' tre khổng lồ được bện xoắn như sợi dây thừng vặn chặt, rối tung đến mức trông như một cơn ác mộng về kiến trúc. Nó không tuân theo bất kỳ quy tắc thẩm mỹ nào, vừa lộn xộn vừa rối rắm, và hoàn toàn không giống bất kỳ định nghĩa nào của cây tre lễ Thất tịch. Trông nó đáng sợ hơn là một dây leo ma quái, và nếu không biết trước, ai cũng có thể nhầm nó là đạo cụ trong một buổi tế lễ cổ xưa nào đó.

Ryuuhou đứng bên cạnh ông anh trai, nhìn cấu trúc tre kỳ quặc kia với vẻ mặt vừa kinh ngạc vừa buồn cười. Cô bé lắc đầu, giọng đầy hài hước: ``Noel-san, Nene-san, hai chị định dựng tổ quạ hay là cây tre Thất tịch vậy? Trông nó giống hệt tác phẩm trừu tượng của một họa sĩ đang lên cơn sốt ấy.''

Mel há hốc mồm kinh ngạc, Flare lùi lại vài bước không dám tin, còn Kuon thì nheo mắt lại như đang cố giải một phương trình vô nghiệm. Cả ba người cùng hét lên đồng thanh: ``\textbf{Cái gì thế này!?}''

Kuon chỉ tay vào ``tác phẩm'' dị hợm kia, giọng không giấu được sự bối rối: ``Cái này\dots\ nhìn chả giống tre tí nào luôn ấy!''

Ngược lại, nàng Hải cẩu vẫn hớn hở chẳng chút nghi ngờ, giơ một tờ đơn từ đâu ra và tuyên bố tự hào, như thể vừa công bố một chiến thắng vĩ đại: ``Giờ chỉ cần gửi đi cuộc thi sáng tác tre thôi! Kiểu gì bọn mình cũng chẳng thắng đâu!''

Ryuuhou nghe vậy, suýt nữa thì phun nước bọt ra vì cười quá lớn: ``Nene-san có chắc là muốn gửi cái này đi dự thi không ạ? Em sợ ban giám khảo sẽ bị sốc vì `nghệ thuật' quá đấy!''

Ba người còn lại cùng lúc ngã ngửa xuống sàn như nhân vật hoạt hình. Nàng Dơi chống tay ngồi dậy trước, bối rối hỏi: ``Cuộc thi nào mà kỳ lạ thế kia!?''

Cậu Tổng xua tay, thở dài như một ông cụ non: ``Anh cũng chả hiểu, nhưng nhìn cái này là hỏng rồi. Không được đâu.''

Nàng Yêu tinh vàng nhướng mày, khoanh tay và nói thẳng với giọng quyết đoán: ``Đập đi, làm lại từ đầu!''

Ryuuhou cũng gật đầu đồng tình, đứng dậy phủi bụi trên quần: ``Em đồng ý với Flare-san. Cái này mà treo lên, chắc mấy vị thần trên trời cũng phải bỏ chạy mất.''

Không chần chừ một giây, Noel gật đầu rất nghiêm chỉnh: ``Rõ!'' rồi lập tức lao vào phá vỡ cái ``dây tre'' nhì nhằng kia, hai tay cô nàng cuồn cuộn sức mạnh xé toạc từng đoạn tre một cách dễ dàng.

Vài phút sau, nhờ sức mạnh cơ bắp phi thường, Noel đã tháo dỡ toàn bộ ``tác phẩm'' cũ. Lần này, cùng với Nene, hai người dựng lại thành một ống trúc khổng lồ đặt ngang trên hai chiếc kệ. Nó trông không quá phức tạp, nhưng gọn gàng và rõ ràng hơn rất nhiều. Nhìn vào, ai cũng cảm thấy\dots\ quen quen.

\img{9-6f}

``Khoan đã\dots'' Kuon nghiêng đầu, nheo mắt nhìn chằm chằm vào những thanh tre. ``Anh thấy cái này sao quen thế nhỉ?''

Mel, Flare và Kuon cùng lúc nhìn chằm chằm vào bộ ống tre. Không khí chững lại vài giây, rồi như thể tất cả vừa sực nhớ ra một điều cực kỳ quan trọng\dots

Cả Noel và Nene cùng lúc trợn tròn mắt, hét lên với giọng đầy hoảng hốt: ``\textbf{Thôi chết rồi!}''

Nàng Hải cẩu đập nhẹ tay vào má, mặt méo xệch đi: ``Bọn mình quên mua mì somen rồi!!'' Ryuuhou nhìn Nene với vẻ mặt thương hại: ``Nene-san, chị định dựng ống mì mà lại quên mua mì á? Hay là chị muốn luộc mấy cọng tre này ăn thay thế?''

Ryuuhou cũng nhìn mãi vào ống tre, rồi đột nhiên bật cười: ``Onii, em nhớ rồi! Cái này giống hệt món đồ ống trúc từ năm ngoái! Năm mà Korone-san với Okayu-san dùng tôm của mình để đua đúng không!?''

Kuon cũng giật mình, rồi vỗ trán như vừa nhớ ra chuyện cũ: ``Ờ đúng rồi! Năm ngoái anh ăn mì somen với nhóm Fubuki-chan\footnote{Xem lại {\flaregothic ÉPISODE 69}, quyển số Năm.}! Cái này chính là ống mì trôi nước mà-- ơ này!''

Cậu chưa kịp nói hết câu thì Mel đã bước đến, tay xách một túi ni-lông từ lúc nào không hay, trên mặt nở nụ cười đầy bí ẩn. ``Chị có mua mì somen đây\~{}''

Kuon lập tức nghi ngờ, nheo mắt nhìn Mel với vẻ dò xét: ``Em không mua nó với giá cắt cổ chứ?''

Nàng Dơi quay mặt đi, cười trừ một cách đáng ngờ: ``À thì\dots\ không đắt lắm đâu ạ\dots\ chắc vậy.''

Ryuuhou nhìn cô, rồi nhìn túi mì, giọng đầy nghi hoặc: ``Mel-san, em có cảm giác chị đang giấu điều gì đó. Cái túi mì này có vẻ hơi đắt tiền quá thì phải.''

Ở một góc khác, Flare cũng chẳng mất nhiều thời gian để phản ứng, liền thốt lên: ``Làm lại lần nữa!''

``Lần nữa hả!?'' Noel và Nene cùng nói, vẻ mặt vừa mệt mỏi vừa cam chịu. Nhưng cuối cùng hai người vẫn đồng ý. Lần thứ ba bắt đầu, Noel lao vào dựng lại từ đầu với một quyết tâm mới.

Lần này, Noel dựng thành một cái giá tre hình trái tim, trông vừa kỳ công vừa\dots\ chẳng liên quan gì đến ngày lễ Thất tịch. Nene còn treo cả đèn lồng giấy lên, khiến toàn bộ cấu trúc trông như một khu vực chụp ảnh dành cho ngày Valentine.

Mel nhìn một hồi, rồi buông một câu đầy thực tế, tay chống hông: ``Giờ thì nó giống quầy chụp ảnh cưới hơn là tre lễ hội.''

Ryuuhou nhìn cấu trúc kia, rồi quay sang Kuon, giọng đầy hài hước: ``Onii này, nếu anh định lấy vợ thì đây có khi là địa điểm chụp ảnh đẹp đấy. Tiếc là nó được làm từ tre thôi.''

Kuon thở dài, nhưng Nene vẫn rất tự tin: ``Thế càng tốt! Đây là tre của tình yêu mà!''

Chưa dừng lại ở đó, lần thứ tư, Noel dựng thành một mô hình xoắn ốc ba tầng, các bậc được lót giấy màu và còn gắn cả bánh xe có đèn nhấp nháy, trông như một công trình kỳ lạ bước ra từ một bộ phim khoa học viễn tưởng.

Flare nhận xét: ``\dots Cái này thì nhìn như tàu lượn siêu tốc bằng tre\dots'' giọng cô vừa hoài nghi, nhưng cũng pha chút thán phục.

``Thế cũng còn tốt hơn cái đống ADN đột biến lúc đầu,'' Kuon thở phào, đưa tay xoa cổ đang mỏi. ``Nhưng mà\dots''

Ryuuhou nhìn cấu trúc xoắn ốc, rồi cũng nhận xét: ``Em thấy cái này có thể làm đồ chơi cho trẻ em được đấy. Nhưng mà liên quan gì đến Thất tịch đâu.''

``Cứ giữ lại hết mấy cái bản này đi,'' Mel đề xuất, giọng đầy nhiệt tình như một nhà sưu tập nghệ thuật đương đại: ``Sau này tổ chức triển lãm nghệ thuật tre vị lai cũng được.''

Cả căn phòng im lặng vài giây, rồi tất cả cùng bật cười, tiếng cười vang vọng khắp phòng như xua tan đi bao mệt mỏi sau mấy lần dựng lại.

``\dots Không nói đùa đâu, anh nghĩ chúng ta vẫn phải làm lại lần nữa,'' Kuon nói, nụ cười trên môi dần tắt, rồi thở ra một hơi dài đầy bất lực.

Ryuuhou vỗ vai anh trai, giọng vừa động viên vừa pha chút trêu chọc: ``Cố lên Onii, đến lần thứ năm thì chắc chắn sẽ thành công. Em sẽ đứng đây ủng hộ Onii hết mình!''

\ct

Cuối cùng, cả sáu người cùng gắng sức dựng lại được một cây tre ngắn gọn hơn hẳn, đặt nằm ngang trên sàn phòng.

\img{9-6g}

Đó là một cây tre thật thà, đơn giản, chẳng có dây leo quấn quýt, chẳng có hình trái tim lệch lạc, cũng chẳng có mô hình tàu lượn kỳ cục. Một cây tre đúng chuẩn: những thanh tre được xếp song song, vững chãi, không tô vẽ cầu kỳ nào, hoàn toàn giống với một cây tre lễ Thất tịch phải thế.

Và rõ ràng, đây chính là thứ mà Noel và Nene đáng lẽ phải làm ngay từ lúc đầu.

Hai cô nàng ngắm nhìn thành quả mới với đôi mắt long lanh, như thể vừa chạm tay vào kiệt tác nghệ thuật đỉnh cao nhất thế giới. Nene cười toe toét, vỗ tay reo lên: ``Qua tay Danchou thì mọi chuyện dễ như ăn bánh thôi!'', đôi mắt cô nàng sáng rực như ngàn vì sao.

``Yay\~{}!'' Noel nối ngay lời, giơ cao nắm tay lên trời như một võ sĩ vừa đăng quang vô địch.

Phía kia, Mel và Flare nằm bẹp trên sàn, tay chân rũ rượi như vừa trải qua cả một khóa huấn luyện sinh tồn đầy khắc nghiệt, sau khi quẩn quanh cùng hai cô nàng lạc lối kia. Cả hai vừa thở hổn hển vì kiệt sức, vừa than thở: ``Mệt chết mất\~{}\dots''

Ryuuhou cũng đứng bên cạnh, thở nhẹ ra một hơi nhưng vẫn giữ được nụ cười trên môi: ``Cuối cùng cũng xong. Em cứ tưởng các chị sẽ mày mò đến tận đêm nay mất.''

Còn Kuon thì khoanh tay, đứng im lặng ngắm nhìn\dots\ rồi lắc đầu ngao ngán, giọng đầy bất lực: ``Tại sao hai đứa không dựng như thế này ngay từ đầu đi, lại còn dày công sáng tạo mấy thứ vớ vẩn kia làm chi nữa\dots''

Đúng lúc đó, Nene vui vẻ giơ cao cuộn giấy màu ngũ sắc lên, giọng phấn khích gọi mọi người: ``Xong rồi nào mọi người ơi, đến lúc viết điều ước thôi\~{}!''

Ngay lập tức, ba cô gái còn lại reo lên hào hứng, xúm lại lấy giấy, mỗi người một màu, cầm bút hí hoáy bắt đầu viết. Kuon chỉ nhướng mày, ngó đồng hồ rồi thở dài: ``Làm vậy có ăn thua gì không đâu? Trễ gần hai mươi ngày rồi đó.''

Ryuuhou nhìn anh trai với ánh mắt tinh nghịch: ``Onii ơi, đã muộn thế này rồi thì trễ thêm mấy ngày nữa cũng có sao đâu. Ít nhất bọn em vẫn còn tre để treo giấy ước mà.''

Nhưng khi nhìn thấy cả bốn cô gái ấy tươi cười, cần mẫn viết ước nguyện của mình trên những dải giấy đủ sắc màu: Nene vẽ một ngôi sao lấp lánh, Noel tô một bông hoa rực rỡ, Mel ghi một dòng chữ nhỏ xíu, Flare viết nắn nót từng nét. Kuon cũng chẳng còn nói gì thêm, cậu chỉ khẽ nhún vai, nhỏ giọng thì thầm với chính mình: ``Thôi thì kệ vậy, chắc mình cũng nên viết một tấm nhỉ?'' rồi cậu cũng bước tới, cầm lấy một dải giấy màu xanh lam, và bắt đầu viết điều ước của riêng mình, dù thầm biết rằng những điều ước dễ thành hiện thực nhất, chính là những điều giản đơn nhất đời thường.

\ct

Theo như Nene đã gợi ý, mỗi thành viên trong nhóm đều viết một điều ước, không chỉ dành cho riêng mình mà còn là những mong muốn tốt đẹp dành cho những người còn lại.

``Mong rằng mình sẽ luôn được theo đuổi đam mê của bản thân.''

``Mong rằng tất cả mọi người sẽ ngày càng trưởng thành và tiến bộ hơn.''

``Mong rằng những kỷ niệm của ngày hôm nay sẽ không bao giờ phai nhạt.''

Sau đó, cả nhóm cùng nhau treo tất cả những dải giấy ước lên cây tre, tay nhẹ nhàng chạm vào từng dải màu, như thể đang gửi gắm hết những mong ước chân thành của mình vào từng nếp giấy nhỏ.

\begin{figure}[H]
  \centering
  \begin{subfigure}[t]{0.45\textwidth}
    \centering
    \includegraphics[width=8cm]{../../../../Image/Book9/9-6h1.png}
  \end{subfigure}
  \quad
  \begin{subfigure}[t]{0.45\textwidth}
    \centering
    \includegraphics[width=8cm]{../../../../Image/Book9/9-6h2.png}
  \end{subfigure}
  \quad
  \begin{subfigure}[t]{0.45\textwidth}
    \centering
    \includegraphics[width=8cm]{../../../../Image/Book9/9-6h3.png}
  \end{subfigure}
  \caption{Các thành viên trong nhóm viết những lời chúc rồi dán lên cây tre.}
\end{figure}

Bầu không khí lúc này trở nên yên ả, dường như chạm đến cảm xúc thiêng liêng của một khoảnh khắc thật đặc biệt. Ánh đèn vàng ấm áp từ những chiếc đèn lồng nhỏ rọi lên những dải giấy đủ màu sắc, tạo nên một khung cảnh đẹp tựa như trong một bức tranh\dots

\dots Cho đến khi Nene bất chợt lên tiếng, giơ tờ kịch bản trong tay lên với đôi mắt tròn rộng vì bối rối: ``Dải ngân hà\dots\ À, ừm\dots\ cái chữ này đọc sao vậy nhỉ? Khó quá đi!''

\img{9-6i}

Nghe cô nàng Hải cẩu nói vậy, Kuon và Flare đồng loạt ngả ngửa ra phía sau, vừa ôm trán vừa bất lực. Kuon lắc đầu ngao ngán: ``Trời ơi\dots\ Cái không khí vừa xây dựng được mấy giây mà bị em phá tan hết rồi!''

Ryuuhou cũng bật cười lớn, tay che miệng để ngăn tiếng cười vang lên: ``Nene-san, chị đã làm hỏng hết cái không khí thiêng liêng mà bọn em vừa cố gắng dựng lên rồi. Nhưng thôi, đó cũng chính là điểm đáng yêu của chị mà.''

Nhưng câu chuyện vẫn chưa dừng lại ở đó.

Bỗng Mel đứng dậy, bước tới phía cửa sổ rồi mở hé cánh cửa ra để ngắm bầu trời đêm. Gió nhẹ thổi vào, làm bay mấy dải giấy trên cây tre vẫy vùng. Cô nhíu mày quan sát rồi nói: ``\dots Ừm, bên ngoài trời mây bao phủ hết cả rồi kìa.''

Kuon cũng bước lại gần, ngó ra ngoài rồi gật gù như thể đang nhớ về một kỷ niệm cũ: ``Giống hệt như chuyện Thất Tịch hai năm trước\footnote{Xem lại {\flaregothic ÉPISODE 9}, quyển số Một.} vậy. Hồi đó anh có nghe nói là Akutan đã giúp Hikoboshi gặp được người yêu của mình đúng không?''

Ryuuhou nhìn lên bầu trời đầy mây, giọng có chút lo lắng: ``Nếu trời cứ mây mù thế này thì Ngưu Lang và Chức Nữ có gặp được nhau không nhỉ? Hay là hai người lại phải đợi thêm một năm nữa mới có thể gặp mặt?''

Nene cũng bắt đầu thấy lo lắng, vừa lắc nhẹ tay Noel vừa nói với giọng bất an: ``Thế thì điều ước của bọn mình sẽ không thành hiện thực mất!''

``Mấy đứa có nghĩ rằng những điều ước lại quan trọng đến thế sao?'' Kuon nhướng mày, giọng như một người đã trải qua nhiều chuyện. ``Nếu mong ước này không thành được vào dịp này thì cũng sẽ có dịp khác mà thôi.''

Ryuuhou nhìn anh trai, gật đầu đồng tình với lời nói của cậu: ``Onii nói đúng lắm. Điều quan trọng nhất là mình đã có những khoảnh khắc vui vẻ bên nhau, chứ không phải chỉ là những điều ước viết trên mảnh giấy.''

Nhưng Noel, với đôi mắt sáng rực như một chiến binh vừa bước ra chiến trường, giơ tay chỉ thẳng lên bầu trời và nói với đầy quyết tâm: ``Cứ để Danchou lo mọi thứ!''

``Ê\dots\ Ê! Em định làm gì thế?'' cậu Tổng vội vàng ngăn cản, nhưng đã quá muộn. Cô em gái cũng trợn tròn mắt, ngay lập tức hiểu Noel định làm gì: ``Noel-san! Khoan đã, cái cửa sổ đấy--!''

Không ai có thể kịp ngăn lại hành động của cô nàng.

Không cần suy nghĩ, không cần đắn đo, Noel nhấc bổng cả cây tre lên và ném thẳng ra ngoài cửa sổ, làm vỡ tung cả khung kính. Cây tre bay vút lên không trung như một tên lửa vừa rời bệ phóng, xé toạc màn mây dày đặc, vẽ nên một đường cong hoàn hảo hướng về bầu trời đêm.

\img{9-6j}

Nhìn cây tre được quăng lên cao, Nene phấn khích hú hét, đôi mắt cô nàng sáng rực như một đứa trẻ vừa được xem pháo hoa: ``\textbf{Tuyệt vời quá đi!!}''

Còn Mel và Flare thì tròn mắt nhìn cảnh tượng vừa xảy ra, miệng há hốc không thể tin được: ``C-Cái gì vậy trời!???''

Ryuuhou đưa tay lên trán, lắc đầu cười khổ: ``Noel-san, chị đúng là chẳng bao giờ biết chữ `dừng lại' viết thế nào nhỉ. Thế là bây giờ Onii lại phải mất công lắp thêm cửa kính mới rồi đấy.''

Kuon thì chỉ biết ôm trán, tặc lưỡi chẳng còn muốn nói gì thêm, chỉ lạnh lùng thốt ra mấy từ với giọng đầy cam chịu: ``\dots Trừ 1500 yên tiền sửa cửa sổ nhé, Hiệp sĩ PON.''

\begin{center}
  \uline{Mặt Trăng.}
\end{center}

Trên bề mặt lấp lánh ánh bạc, vị thần Fubuoshi đang ngồi thảnh thơi bên con trâu của mình, tay áo bay phấp phới như một vị tiên cổ điển. Dù trông hoàn toàn giống Fubuki, gần chỗ ngồi của ông có treo một tấm biển gỗ viết mấy chữ rất to:

\txtbx{``Fubuki-chan là ai? Ta là Fubuoshi. Mấy người nhầm ta với ai đó sao!?''}

Vị Thần Mặt Trăng đang nhắm mắt suy tư sâu xa về vũ trụ và những mối kết nối tâm linh, hay nói đúng hơn là đang nhớ nhung đến Chức Nữ, thì\dots

*VỤT!*

\dots ông giật mình bởi một vật thể bất ngờ lao tới từ hạ giới.

\img{9-6k}

Một cây tre bay thẳng về phía ông như một thiên thạch xuyên không, cắm chặt vào đám mây nơi ông đang ngự. Lực va chạm làm rung chuyển cả mặt phẳng mây, khiến con trâu của ông cũng phải giật mình lùi lại vài bước.

Fubuoshi chẳng chút nghi ngờ, cưỡi lên con trâu của mình, lóc cóc đi tới gần cây tre rồi cúi xuống nhìn những dải giấy lấp lánh vẫn còn treo lủng lẳng trên đó. Vì tò mò, ông cầm lấy một tờ giấy màu cam lên đọc. Trên giấy là một dòng chữ nguệch ngoạc, trông như được viết vội vàng: ``Hmm\dots\ Đây là điều ước của\dots\ Suzumomo?''

Điều ước viết trên tờ giấy đó vô cùng đơn giản, đến nỗi người đọc phải nhắm chặt mắt lại để chắc chắn mình không nhìn lầm: ``Nene muốn ăn somen!''

Vị Thần nhìn dải giấy, im lặng vài giây rồi cụp tai nhăn mặt khó hiểu nói: ``\dots Thế thì tự đi ra quán ăn đi?'' giọng ông đầy bất lực, như thể đang nói chuyện với một đứa trẻ con ngây thơ.

Và đó mới chỉ là dải giấy đầu tiên trong số sáu điều ước đang chờ ông giải mã. Ông lật sang những dải giấy khác, trong đó có một dải giấy màu xanh lá với nét chữ thanh mảnh, viết rằng: ``Mong mọi người luôn vui vẻ và mạnh khỏe.'' Không có tên người viết, nhưng nếu nhìn kỹ, có thể thấy một hình vẽ nhỏ ở góc giấy: một cây đàn guitar và một nốt nhạc.

%==========================================TRUYỆN PHỤ=========================================
\omk

Fubuoshi lần lượt đọc những mảnh giấy ước còn lại từ cây tre vừa bị quăng lên trời. Ông nhặt lên một tờ với chữ viết đều đặn, khẽ tự nhủ: ``Hừm\dots `Noel muốn ăn thật nhiều gyudon'\dots?'' Vị Thần Cáo trắng nhíu mày, liếc qua tờ giấy của Nene rồi thở dài một cách mệt mỏi: ``Lại thêm một người ước đồ ăn nữa\dots''

Tiếp theo, ông lật qua mảnh giấy của Mel và gật gù một cách hài lòng, ánh mắt ông rạng rỡ lên chút ít: ``Mel mong mọi người ở \hll\ luôn mạnh khỏe và có thể theo đuổi đam mê của mình thật tốt.''

Ngay sau đó là dải giấy của Flare, với nét chữ mềm mại và chỉn chu: ``Con ước rằng chúng ta sẽ luôn được cháy hết mình trên sân khấu cùng những người mà mình trân quý.''

Vị Thần nhìn hai lời ước đó và nở một nụ cười ấm áp, ánh mắt tràn đầy sự yêu mến: ``Cuối cùng cũng có mấy đứa nghiêm túc. Hai đứa kia ở dưới chắc chắn là tấm gương tốt cho mọi người.''

Nhưng rồi ông dừng lại khi cầm lên tờ giấy thứ năm, viết trên dải giấy màu xanh đậm, nét chữ hơi xiên xẹo, trông như được viết vội vàng giữa lúc suy tư: ``Tôi mong nhóm này là một nhóm idol, chứ đừng trở thành những gã hề rạp xiếc.''

Một khoảng lặng ngắn ngủi. Fubuoshi ngước nhìn xuống hạ giới, gương mặt pha lẫn giữa bật cười và chút xót xa, ông đưa tay xoa cằm: ``Cậu bé này\dots\ nói đúng mà cũng thật là làm người ta đau lòng.''

Chưa kịp nghĩ ngợi thêm, một mảnh giấy khác bất ngờ bay tới, đập thẳng vào mặt ông. Mở ra xem, ông đọc được một điều ước nữa với nét chữ giống hệt dải xanh đậm vừa rồi, nhưng lần này giọng điệu lại có phần trêu chọc hơn: ``Và mong ông Thần ban cho bọn con mấy cái cửa sổ chống đạn nhé.''

Fubuoshi im lặng đúng ba giây. Rồi ông bật cười thành tiếng, vừa lắc đầu vừa lẩm bẩm: ``Đến cả những điều ước phụ cũng không quên châm chọc nhau, đúng là\dots''

Cuối cùng, ông cầm lên dải giấy cuối cùng - màu xanh lá cây, với nét chữ thanh thoát nhưng đầy tự tin, trông rất giống chữ của một người đã quá quen với việc đứng trước đám đông. Trên đó, một dòng chữ nhẹ nhàng mà vô cùng sâu sắc: ``Em mong mọi người sẽ luôn được sống trong những khoảnh khắc đáng nhớ như hôm nay. Và em mong ban nhạc của em - những người bạn thân nhất của em - sẽ luôn ở bên cạnh em, dù có chuyện gì đi nữa.''

Fubuoshi nhìn dải giấy, khẽ mỉm cười một cách trìu mến. Ông chẳng nói thêm lời nào, chỉ nhẹ nhàng cất nó vào một vị trí đặc biệt trong tay áo, như thể đó là điều ước quý giá nhất mà ông từng nhận được.

Không chần chừ thêm gì nữa, ông cưỡi con trâu của mình bay thẳng xuống trần gian.

\ct

Tại cái cửa sổ vừa bị vỡ lúc trước, cả sáu người đang xúm lại đánh giá mức độ hư hại. Cậu Tổng vừa mò tìm số điện thoại của thợ lắp kính thì\dots

*BÙM!*

Một luồng sáng trắng chói chang đổ thẳng xuống sân. Từ trên cao, một người cưỡi trên lưng trâu đáp xuống, tay vẫn còn ôm đống dải giấy ước, còn trên trán thì hằn rõ vết đỏ do bị mảnh giấy đập vào.

\img{9-6l}

``Cũng đã hai năm trôi qua kể từ lần cuối chúng ta gặp mặt rồi nhỉ,'' Fubuoshi cất lời, ánh mắt hiền hậu hướng về phía Kuon. ``Có vẻ như cậu đã cầu xin hơi nhiều rồi đấy'' ông cười, giọng nói ẩn chứa nhiều ý nghĩa.

Kuon cúi gập người, giọng pha lẫn chút ngượng ngùng: ``Thưa ngài Fubuoshi-sama! Dạ, cũng không giấu gì ngài, cháu có lẽ đã đòi hỏi hơi\dots\ quá mức ạ.'' Nene thì nhìn bóng người đang từ từ hạ xuống kia, liền buột miệng một cách ngây thơ, giọng vẫn tràn đầy ngạc nhiên: ``Kia chẳng phải là Fubuki-senpai sao?''

Nghe vậy, vị thần giật mình, rồi vội vàng phản đối, giọng đầy bực bội: ``Không phải! Ta không phải Fubuki! Tên ta là Fubuoshi! FU-BU-O-SHI! Các ngươi đang nhầm lẫn ta với ai thế!?'' ông vung tay lia lịa lên trời, trông chẳng khác nào một thầy giáo đang cố gắng dạy đám học sinh lớp một nhớ bài.

Ryuuhou đứng sát bên Kuon, nhướng mày quan sát vị thần, rồi quay sang anh trai mình, thì thầm nhỏ: ``Onii, ông ấy giống y hệt Fubuki-san. Em không tin là họ không phải cùng một người đâu.''

``Thì đúng là em ấy còn gì,'' ông anh nhún vai đáp lại, ``em ấy chỉ pha trò cho vui thôi.''

Noel nghiêng đầu, trả lời ngắn gọn với giọng trong sáng: ``Thì có đâu, tại trông ông ấy quá giống Fubuki-senpai mà\dots''

Flare nhíu mắt lại quan sát kỹ, rồi cũng gật gù đồng tình: ``Nhìn kỹ lại thì đúng là rất giống thật ạ\dots''

Ryuuhou lại thêm vào, giọng cứ như một nhà phê bình chuyên nghiệp: ``Cả cái kiểu cau có mày cũng giống hệt. Chắc chắn ông ấy là họ hàng xa của Fubuki-san rồi.''

Vị thần cáo trắng nghe đến đây chỉ còn biết ngửa cổ than trời, mặt đầy vẻ bất lực, hai tay tung lên trời: ``Tại sao ai cũng nói vậy chứ\~{}??''

Cả sáu người cùng bật cười trước phản ứng giận dỗi của ông vì sự trùng hợp kỳ lạ này. Fubuoshi thở dài một hơi, rồi hắng giọng cố trở nên nghiêm túc, gắng gượng lấy lại phong thái của một vị thần: ``Ta đã đọc hết tất cả những điều ước của các ngươi rồi. Flare, Mel và Kuon, lời nguyện của ba ngươi\dots\ nghe rất thực tế và đáng trân trọng.''

Rồi ông quay sang Nene và Noel, liếc dài một cái với ánh mắt không còn biết nói gì: ``Còn hai ngươi thì\dots\ ước cái kiểu gì thế trời. Ai lại đi ước được ăn cơm bò và mì somen chứ?''

Ryuuhou cười khúc khích, nhìn Nene và Noel với vẻ đáng thương: ``Nene-san, Noel-san, hai chị có thể ước một cái gì đó lớn lao hơn không? Em tưởng hai chị sẽ ước được đứng trên sân khấu lớn hay gì đó tương tự chứ.''

Cặp đôi bị chỉ điểm chỉ biết cười trừ, trong khi ba người còn lại nhìn họ với ánh mắt đầy ``nghi ngờ nhân phẩm''. Kuon lắc đầu, giọng đầy bất lực: ``Anh cũng chưa từng thấy ai ước mấy thứ linh tinh vặt vãnh như hai đứa này bao giờ.''

Fubuoshi lại tiếp tục nói, giọng đã dịu đi nhiều: ``Thôi được rồi. Theo đúng nguyện vọng của các ngươi, ta sẽ ban cho mỗi người một món quà nhỏ.''

Ông búng tay một cái. Chỉ trong nháy mắt, trên tay Nene hiện ra một bát mì somen còn bốc khói nghi ngút, còn Noel thì ôm trọn cả một hộp cơm bò cự khổng lồ. Mel lúc này hơi nhíu mày: ``Ơ, bọn mình vừa mới mua mì mà\dots''

Ryuuhou nhìn hai món quà trước mắt, bật cười lớn: ``Mel-san ơi, chị đã mua mì rồi mà còn ước làm gì nữa. Cái của ông thần đây còn nóng hổi hơn nhiều kia.''

Kuon nhún vai đáp lại: ``Ừ, nhưng thôi kệ. Cái này còn nóng hơn.''

Nene và Noel cùng nhau nói nhanh như chớp, cúi đầu cảm ơn rối rít: ``Cảm ơn chị nhé, Fubuki-senpai!''

``\textbf{Đã bảo là ta không phải Fubuki mà!}'' vị thần gắt lên, giọng vẫn đầy bực bội vì bị gọi nhầm tới lần thứ hai trong ngày.

``Xin lỗi, xin lỗi\~{}'' Nene nói, cố tình đọc sai tên để trêu chọc ông.

Fubuoshi thở dài một lần nữa, rồi lại búng tay thêm cái nữa. Một tấm kính dày cộp, cứng cáp từ từ hiện ra, lắp đúng vào vị trí ô cửa bị vỡ lúc trước. Ánh sáng phản chiếu từ tấm kính mới khiến cả nhóm phải tròn mắt ngạc nhiên.

``Ối, trông nó còn xịn hơn cái cửa kính cũ nhiều,'' Noel gõ nhẹ vào cánh cửa, cảm nhận độ bền của tấm kính, rồi thở phào nhẹ nhõm. Cô quay sang Fubuoshi, giọng đầy biết ơn: ``Cảm ơn ông nhiều ạ!''

Ryuuhou bước tới, đưa tay chạm vào tấm kính, cảm nhận sự chắc chắn của nó, rồi nhận xét: ``Tấm kính này chắc đến nỗi có ai ném búa vào cũng không vỡ đâu. Onii, lần sau anh có thể hoàn toàn yên tâm rồi.''

Kuon vỗ nhẹ vào vai nữ Hiệp sĩ, nở ra một nụ cười hài lòng: ``May nhá. Lần sau nhớ mở cửa trước rồi mới ném cái gì ra ngoài nghe chưa.''

Fubuoshi nhìn cả sáu người đang đứng tụ tập bên cửa sổ, bỗng dưng khựng lại một chút. Ánh mắt ông trở nên dịu dàng hơn bao giờ hết, giống như một người ông đang nhìn đàn cháu nhỏ tung tăng vui đùa.

``Những lời ước năm nay\dots\ có lẽ chẳng có gì to tát, chẳng có gì hào nhoáng, nhưng lại rất chân thật. Cảm ơn các ngươi vì vẫn giữ được trái tim mơ mộng, dù là trong những khoảnh khắc hài hước hay lúc nghiêm túc.''

Nói xong, ông giơ ngón tay cái lên, ánh trăng rọi thẳng vào đôi mắt đang híp lại vì mãn nguyện. Sáu người bên dưới cùng lúc cúi đầu, hô to một tiếng, giọng nói hòa quyện vào nhau như một bản hợp xướng: ``Cảm ơn Hikoboshi-sama!''

Fubuoshi cưỡi con trâu bay lên cao, tay vẫy vẫy tạm biệt mà không nói thêm một lời nào. Bóng ông dần nhỏ lại, hòa tan vào ánh trăng bạc trên bầu trời.

Kuon khẽ nói, mắt vẫn dõi theo bóng vị thần đang bay xa, giọng trầm hẳn đi: ``Hẹn gặp lại\dots\ nếu có thể.''

Ryuuhou đứng cạnh anh, ngước nhìn lên bầu trời, giọng nói nhẹ nhàng như gió: ``Onii, em nghĩ ông ấy vẫn còn nhớ những điều ước của bọn mình đấy. Dù ông ấy có phải là Fubuki-san hay không, thì ông cũng đã giúp chúng ta có một buổi tối thật tuyệt vời.''

Hình ảnh cửa sổ vừa được thay mới phản chiếu ánh trăng dịu nhẹ, những đốm sáng từ thành phố phía xa dường như lấp lánh hơn mọi ngày. Cả nhóm vẫn đứng yên đó, ai cũng cầm trên tay món quà vừa được ban tặng, gương mặt tràn đầy niềm vui.

Mel nghiêng đầu tựa vào vai Flare, thì thầm điều gì đó khiến cả hai cùng cười khúc khích. Nene tranh thủ chụp ảnh tự sướng với bát mì somen của mình, còn Noel thì vừa ăn vừa nói: ``Lần sau chắc ước cả một cái đùi gà quay luôn!'' cô nói, miệng vẫn nhai nhồm nhoàm.

Ryuuhou đứng bên cạnh, tay cầm miếng cơm bò mà Noel chia cho, mỉm cười nhẹ: ``Em chỉ ước là ban nhạc của em sẽ có một buổi diễn đáng nhớ như hôm nay. Không cần cái gì đó quá cao siêu, chỉ cần được ở bên những người mình yêu quý là đủ rồi.''

Kuon nhìn em gái, rồi đưa tay xoa nhẹ đầu cô: ``Rồi sẽ có ngày đó thôi, Ryuuhou. Chỉ cần em luôn tin vào điều đó.''

Cậu ngước nhìn lên bầu trời, nơi đốm sáng trắng của Hikoboshi đang mờ dần đi như một ngôi sao lạc lõng. Cậu không nói thêm lời nào, chỉ đặt tay lên khung cửa mới vừa được lắp, như một lời cảm ơn thầm lặng gửi đến vị thần.

Một làn gió nhẹ lướt qua, mang theo hương vị của đêm hè, của bữa tiệc sau lễ Thất tịch vẫn còn dang dở, của những điều ước nhỏ bé nhưng chứa đựng hết cả tấm chân tình.

\end{document}